% 第4章 4.2.2：(a) 阶跃 H=1_{(0,∞)}（H(0)=0，实心点），广义导数为 delta_0；(b) 支撑缩小的测试 phi_eps(x)=chi(x/eps)，eps=1,0.5,0.25，
% 均有 phi_eps(0)=1、0<=phi_eps<=1、支撑[-0.9eps,0.9eps]。数据由同目录 generate.py 计算。
\begin{tikzpicture}[font=\figurefont\footnotesize]
\pgfplotsset{every axis/.style={width=50mm,height=34mm,axis lines=middle,
  x axis line style={draw=FigureStroke,line width=.6pt,-{Latex[length=1.6mm]}},
  y axis line style={draw=FigureStroke,line width=.6pt,-{Latex[length=1.6mm]}},
  tick style={draw=FigureStroke},tick label style={font=\figurefont\scriptsize,text=FigureInk},
  xlabel={$x$},xlabel style={anchor=west,font=\figurefont\scriptsize},
  title style={font=\figurefont\footnotesize,yshift=1mm},xmin=-1.2,xmax=1.3,ymin=-0.15,ymax=1.3}}
\begin{axis}[name=A,xtick={-1,1},ytick={1},title={(a) 阶跃$H$}]
\addplot[draw=FigureBlue,line width=1.1pt,domain=-1.1:0] {0};
\addplot[draw=FigureBlue,line width=1.1pt,domain=0:1.1] {1};
\addplot[only marks,mark=o,mark size=1.4pt,draw=FigureBlue,fill=white] coordinates {(0,1)};
\addplot[only marks,mark=*,mark size=1.4pt,FigureBlue] coordinates {(0,0)};
\end{axis}
\begin{axis}[at={(A.east)},anchor=west,xshift=12mm,axis y line=left,xtick={-0.9,0.9},ytick={1},title={(b) 测试$\varphi_\varepsilon(x)=\chi(x/\varepsilon)$},
  legend style={draw=none,fill=none,font=\figurefont\scriptsize,at={(1.03,1)},anchor=north west,cells={anchor=west}}]
\def\dat{figures/02-abstract-mathematics/03-function-spaces/tikz/c4-step-delta/tests.dat}
\addplot[draw=FigureYellow,line width=1.1pt,dash dot] table[x=x,y=e1] {\dat};
\addplot[draw=FigureGreen,line width=1.1pt,dashed] table[x=x,y=e2] {\dat};
\addplot[draw=FigureBlue,line width=1.1pt] table[x=x,y=e4] {\dat};
\legend{$\varepsilon=1$,$\varepsilon=0.5$,$\varepsilon=0.25$}
\end{axis}
\end{tikzpicture}
